\chapter*{LEMBAR PENGESAHAN}
\addcontentsline{toc}{chapter}{LEMBAR PENGESAHAN}

\begin{center}
  \vspace{2cm}
  {\large\bfseries Pengembangan GraphRAG dengan Pembobotan Graf Pengetahuan Berbasis Penilaian Pengguna untuk Mengurangi Halusinasi pada \emph{Large Language Model}}\\
  \vspace{2cm}

  {\Large \textbf{TESIS}}\\

  \vspace{1.5cm}

  {\large Oleh}\\[0.3cm]
  \textbf{
    {\large Dartika Anie Marian}\\
    {\large 23524319}
  }\\

  \vspace{0.5cm}

  {\large Program Studi Magister Informatika}\\
  {\large Sekolah Teknik Elektro dan Informatika}\\
  {\large Institut Teknologi Bandung}\\

  \vspace{1.5cm}

  Tesis ini telah disetujui dan disahkan\\
  di Bandung, pada tanggal \today\\[1cm]

  % ==========================================
  % Versi 1 pembimbing (default)
  % ==========================================
  Pembimbing \\[3cm]
  Dr.Ir. Arry Akhmad Arman, MT. \\[0.2cm]
  NIP. 196504141991021001
  % ==========================================
\end{center}

\vspace{1cm}
\noindent

% ==========================================
% Jika ada 2 pembimbing TA, uncomment dan edit
% tabular di bawah ini. Kemudian, comment out atau hapus
% bagian versi 1 pembimbing di atas.
% ==========================================